\section{A refined rational motif certificate}
\label{numerical:section}

This section uses the upper motif certificate from~\cite{UzayTables}.
The numerical data are pinned separately from the later revision used
for the existence proof. The unchanged tables and their original
all-state drift decomposition are inputs; the refinement consists of
a correlated upper weight and a constant correction per core.

\subsection{The potential and its grouped upper weight}

Root the real core tree at a core vertex. A core has type $(L,k,h)$,
where $L$ is its leaf count, $k$ its number of core children, and
$h\in\{0,1\}$ indicates a core parent. The original potential is
\begin{align}\label{num:potential}
 \Phi(T)={}&\sum_v f(v)+\sum_{p\to w}g(p,w)
       +\sum_{p\to u\to w}h_V(p,u,w)\notag\\
       &+\sum_u\sum_{\{w_1,w_2\}\subset\operatorname{ch}_K(u)}
                       h_C(u;w_1,w_2).
\end{align}
All tables have denominator $10^{10}$. The vertex table saturates at
$L_F=800$ and $K_F=150$. For an endpoint of full degree $D\ge24$,
the interpolation measure is
\[
 \pi_{D,k,h}=(1-24/D)\delta_{\mathrm{INF}}
       +(24/D)\delta_{T_1(\min(k,13),h)}.
\]
The vertical middle has at most eight core children and degree below
$16$; the cherry middle has at most six core children and degree below
$12$. Terms outside these middle domains vanish. Endpoint evaluations
are multilinear in the displayed measures. The exact arrays and
projection rules are generated by the pinned \texttt{cert\_build.c}.

Assign an edge and a vertical path to their bottom core, and a cherry
to its middle core. Let $W(v)$ be the resulting contribution at $v$,
so that $\Phi=\sum_vW(v)$. If $\mathcal P,\mathcal W$ are the original
finite parent and child endpoint sets, put
\begin{align*}
 R(w)=\max\bigl\{0,&\ \max_{p\in\mathcal P}g[p,w],\\
 &\max_{u\in H_u,\ p\in\mathcal P}
                 (g[u,w]+h_V[p,u,w])\bigr\},\\
 H^+(u)&=\max\{0,\max_{w_1,w_2\in\mathcal W}h_C[u,w_1,w_2]\}.
\end{align*}
Here $H_u$ is the supported vertical-middle set and $C_u$ below is
the supported cherry-middle set. Define
\begin{equation}\label{num:Q}
 Q(v)=f(v)+h\sum_w\pi_v(w)R(w)
       +\mathbf1_{\{v\in C_u\}}\binom{k}{2}H^+(v).
\end{equation}
The edge and vertical terms in $R$ use the same parent. Positivity of
the interpolation weights therefore gives the pointwise bound
$W(v)\le Q(v)$. Exact evaluation gives
\begin{equation}\label{num:comparison}
 m C\le\Phi\le\overline A C,
 \qquad m=\frac{2989767809}{10^{10}},
 \qquad \overline A=\frac{8424600447}{5000000000}.
\end{equation}
The lower bound is part of the original certificate. The upper bound
uses $\sup_vQ(v)=\overline A$.

Write $\nu=1133/2000$, and let $B_f(v;T)$ be the original assigned
vertex drift. Its vertex reserve is
$\delta(v;T)=\nu f(v)-B_f(v;T)$. The verified nonnegative edge,
vertical, and cherry residuals imply the all-tree inequality
\begin{equation}\label{num:original-drift}
 \mathcal L\Phi(T)\le\nu\Phi(T)-\sum_{v\in K(T)}\delta(v;T).
\end{equation}
This is the drift decomposition to which the refinement is applied.

\subsection{A constant correction per core}

Set
\begin{gather}\label{num:parameters}
 \beta=\frac{2832242711}{5000000000},\qquad
 \tau=\frac{5294970481}{10^{14}},\qquad
 \Phi_\tau=\Phi+\tau C,\\
 \zeta=\frac{199752266575452391}
              {1200500000000000000000000000}>0.
 \label{num:reserve}
\end{gather}

\begin{theorem}[Corrected full-tree drift]\label{num:drift}
Every finite full tree with at least one core satisfies
\begin{equation}\label{num:corrected-drift}
 \mathcal L\Phi_\tau\le\beta\Phi_\tau-\zeta C.
\end{equation}
Consequently $\mu\le\beta=0.5664485422$.
\end{theorem}
\begin{proof}
Put $b_v=L_v/D_v$. Only promotion increases the core count, so
$\mathcal LC=\sum_v b_v$. Since $\nu>\beta$ and $W(v)\le Q(v)$,
\eqref{num:original-drift} gives
\begin{align*}
 \beta\Phi_\tau-\mathcal L\Phi_\tau
 &\ge\sum_v\left[\delta(v;T)-(\nu-\beta)W(v)
                            +\tau(\beta-b_v)\right]\\
 &\ge\sum_v\left[\delta(v;T)-(\nu-\beta)Q(v)
                            +\tau(\beta-b_v)\right].
\end{align*}
The exact additional finite certificate checks
\begin{equation}\label{num:row-check}
 \delta_i-(\nu-\beta)Q_i+\tau(\beta-b_i)\ge\zeta
\end{equation}
in all $484138$ vertex rows. The minimum is the displayed $\zeta$.

For completeness, the original tail interpolation also proves that
these rows cover all degrees. In the unsaturated outer box only the
incoming leaf-gain intensity varies; the expression is affine in
that intensity, and both endpoints are checked. On the leaf tail the
enclosing parameter rectangle is
\[
 b\in\left[\frac{801}{801+k+h},1\right],\qquad
 t\in\left[0,\frac{24}{801+k+h}\right].
\]
The vertex value is constant there, the assigned drift is bilinear in
$(b,t)$, and $Q$ is affine in $t$. The correction is affine in $b$.
Thus \eqref{num:row-check} is bilinear and follows from its four corner
checks. The post-promotion edge uses child-count label $k+1$ and $Q$
uses the pre-promotion label $k$, at the same value of $t$. The exact
$L=800$ boundary is checked separately.

For overflow $k>150$, the rectangle is
$b\in[0,1]$, $t\in[0,24/(151+h)]$, with the clamped child-count
label $13$. The same four-corner argument applies. Hence
\eqref{num:row-check} holds at every actual vertex, and summation
proves \eqref{num:corrected-drift}.

The lower comparison $(m+\tau)C\le\Phi_\tau$ and the upper comparison
$(\overline A+\tau)C\ge\Phi_\tau$ follow from
\eqref{num:comparison}. Dropping the strictly positive reserve and
applying the resulting drift inequality gives expected growth at most
$e^{\beta t}$ from every fixed finite seed. Equation
\eqref{model:imported-exponent} identifies the original exponent,
so $\mu\le\beta$.
\end{proof}

The exact feasible interval for $\tau$ at this fixed $\beta$ is
\[
 \frac{763278221264673581}{14415721236665000000000}
 \le\tau\le
 \frac{32452312724432471239}{612865341021362000000000}.
\]
The lower constraint is the outer row $(L,k,h)=(143,150,1)$ at its
larger incoming-rate endpoint, with $b=143/294$; the upper constraint
is the analogous row $(795,150,1)$, with $b=795/946$. The chosen
rational $\tau$ is strictly inside this interval. No optimality claim
for this pair is needed.

At the deterministic three-vertex seed,
$\Phi(T_3)=14606753031/10^{10}$ and $C(T_3)=1$. Dividing the generator
inequality by the current population and iterating gives the explicit
finite-size bound
\begin{equation}\label{num:finite-size}
 \mathbb E C(T_n)\le
 \frac{14606753031/10^{10}+\tau}{m+\tau}
 \prod_{j=3}^{n-1}\left(1+\frac\beta j\right),\qquad n\ge3.
\end{equation}

\subsection{Reproduction of the certificate}

The repository supplies \texttt{scripts/verify\_bound.py}, which fetches
the immutable upstream inputs, verifies their hashes, builds the
original rational binary, runs the complete GMP verifier, and checks
the corrected rows in exact arithmetic. The upstream manifest and
the two result files are under \texttt{verification/}. The full run
checks $17329108$ original motif inequalities and $37$ finite trees,
then all $484138$ corrected vertex rows. The tail argument above
extends those finite comparisons to arbitrary degrees.

The corrected row check reads the original arrays; it does not replace
the original edge, vertical, or cherry checks. All decisions are
rational sign comparisons. Floating-point calculations may propose
parameters but are not part of the certificate. Reproduction with the
same verified source is a reproducibility check, rather than an
independently implemented second verifier.
